% GENERATED FILE. Edit canonical lessons, then run npm run generate:books.
% canonical-source-hash: fnv1a64:4914087dde082d28
% canonical-lessons: BN-C213-chul-kata, BN-C213-saja, BN-C213-prostut-hooya, BN-C213-roona-deoya, BN-C213-chora

\chapter{To have a haircut, To dress up, To get ready, To set off, To ride}
\label{ch:bn-to-have-a-haircut-to-dress-up-to-get-ready-to-set-off-to-ride}
\hlchaptermodality{213}

\begin{chapteropening}
\textbf{By the end of this chapter:} \emph{I can say to have a haircut, to dress up, to get ready, to set off and to ride.}

Complete the last lesson of chapter 213: I can say to have a haircut, to dress up, to get ready, to set off and to ride.
\end{chapteropening}

\section[\texorpdfstring{\bn{চুল} \bn{কাটা} (chul kāṭā)}{চুল কাটা (chul kāṭā)}]{\bn{চুল} \bn{কাটা} (chul kāṭā) — to have a haircut}

\label{lesson:BN-C213-chul-kata}



\begin{quote}
Before the new one: say the Bengali for to fight, then the Bengali for to run away.
\end{quote}

\subsection*{You'll want to know: \bn{চুল} \bn{কাটা}}
\textbf{\bn{চুল} \bn{কাটা}} — \emph{chul kāṭā} — \textquotedblleft{}to have a haircut\textquotedblright{}.

\begin{grammarlens}[title={What you've built}]
One more word for this chapter.
\end{grammarlens}

\subsection*{Your turn}
\begin{itemize}
\raggedright
  \item \emph{Say it:} \emph{chul kāṭā}
  \item \emph{Say it:} \emph{chul kāṭā}, once more
  \item \emph{Say it:} say \emph{chul kāṭā}
  \item \emph{Recall:} say the Bengali for to fight, then the Bengali for to run away, then say \emph{chul kāṭā} again
\end{itemize}

\begin{tcolorbox}[breakable,colback=teal!4,colframe=teal!35!black,title={Before you move on}]
What does \bn{চুল} \bn{কাটা} mean? (\textquotedblleft{}To have a haircut\textquotedblright{}.) Say it once more.
\end{tcolorbox}
\section[\texorpdfstring{\bn{সাজা} (sājā)}{সাজা (sājā)}]{\bn{সাজা} (sājā) — to dress up}

\label{lesson:BN-C213-saja}



\begin{quote}
Before the new one: say the Bengali for to run away, then the Bengali for to have a haircut.
\end{quote}

\subsection*{You'll want to know: \bn{সাজা}}
\textbf{\bn{সাজা}} — \emph{sājā} — \textquotedblleft{}to dress up\textquotedblright{}.

\begin{grammarlens}[title={What you've built}]
One more word for this chapter.
\end{grammarlens}

\subsection*{Your turn}
\begin{itemize}
\raggedright
  \item \emph{Say it:} \emph{sājā}
  \item \emph{Say it:} \emph{sājā}, once more
  \item \emph{Say it:} say \emph{sājā}
  \item \emph{Recall:} say the Bengali for to run away, then the Bengali for to have a haircut, then say \emph{sājā} again
\end{itemize}

\begin{tcolorbox}[breakable,colback=teal!4,colframe=teal!35!black,title={Before you move on}]
What does \bn{সাজা} mean? (\textquotedblleft{}To dress up\textquotedblright{}.) Say it once more.
\end{tcolorbox}
\section[\texorpdfstring{\bn{প্রস্তুত} \bn{হওয়া} (prôstut hôoyā)}{প্রস্তুত হওয়া (prôstut hôoyā)}]{\bn{প্রস্তুত} \bn{হওয়া} (prôstut hôoyā) — to get ready}

\label{lesson:BN-C213-prostut-hooya}



\begin{quote}
Before the new one: say the Bengali for to have a haircut, then the Bengali for to dress up.
\end{quote}

\subsection*{You'll want to know: \bn{প্রস্তুত} \bn{হওয়া}}
\textbf{\bn{প্রস্তুত} \bn{হওয়া}} — \emph{prôstut hôoyā} — \textquotedblleft{}to get ready\textquotedblright{}.

\begin{grammarlens}[title={What you've built}]
One more word for this chapter.
\end{grammarlens}

\subsection*{Your turn}
\begin{itemize}
\raggedright
  \item \emph{Say it:} \emph{prôstut hôoyā}
  \item \emph{Say it:} \emph{prôstut hôoyā}, once more
  \item \emph{Say it:} say \emph{prôstut hôoyā}
  \item \emph{Recall:} say the Bengali for to have a haircut, then the Bengali for to dress up, then say \emph{prôstut hôoyā} again
\end{itemize}

\begin{tcolorbox}[breakable,colback=teal!4,colframe=teal!35!black,title={Before you move on}]
What does \bn{প্রস্তুত} \bn{হওয়া} mean? (\textquotedblleft{}To get ready\textquotedblright{}.) Say it once more.
\end{tcolorbox}
\section[\texorpdfstring{\bn{রওনা} \bn{দেওয়া} (rôonā deoyā)}{রওনা দেওয়া (rôonā deoyā)}]{\bn{রওনা} \bn{দেওয়া} (rôonā deoyā) — to set off}

\label{lesson:BN-C213-roona-deoya}



\begin{quote}
Before the new one: say the Bengali for to dress up, then the Bengali for to get ready.
\end{quote}

\subsection*{You'll want to know: \bn{রওনা} \bn{দেওয়া}}
\textbf{\bn{রওনা} \bn{দেওয়া}} — \emph{rôonā deoyā} — \textquotedblleft{}to set off\textquotedblright{}.

\begin{grammarlens}[title={What you've built}]
One more word for this chapter.
\end{grammarlens}

\subsection*{Your turn}
\begin{itemize}
\raggedright
  \item \emph{Say it:} \emph{rôonā deoyā}
  \item \emph{Say it:} \emph{rôonā deoyā}, once more
  \item \emph{Say it:} say \emph{rôonā deoyā}
  \item \emph{Recall:} say the Bengali for to dress up, then the Bengali for to get ready, then say \emph{rôonā deoyā} again
\end{itemize}

\begin{tcolorbox}[breakable,colback=teal!4,colframe=teal!35!black,title={Before you move on}]
What does \bn{রওনা} \bn{দেওয়া} mean? (\textquotedblleft{}To set off\textquotedblright{}.) Say it once more.
\end{tcolorbox}
\section[\texorpdfstring{\bn{চড়া} (chôṛā)}{চড়া (chôṛā)}]{\bn{চড়া} (chôṛā) — to ride, to climb}

\label{lesson:BN-C213-chora}



\begin{quote}
Before the new one: say the Bengali for to get ready, then the Bengali for to set off.
\end{quote}

\subsection*{You'll want to know: \bn{চড়া}}
\textbf{\bn{চড়া}} — \emph{chôṛā} — \textquotedblleft{}to ride, to climb\textquotedblright{}.

\begin{grammarlens}[title={What you've built}]
One more word for this chapter.
\end{grammarlens}

\subsection*{Your turn}
\begin{itemize}
\raggedright
  \item \emph{Say it:} \emph{chôṛā}
  \item \emph{Say it:} \emph{chôṛā}, once more
  \item \emph{Say it:} say \emph{chôṛā}
  \item \emph{Recall:} say the Bengali for to get ready, then the Bengali for to set off, then say \emph{chôṛā} again
\end{itemize}

\begin{tcolorbox}[breakable,colback=teal!4,colframe=teal!35!black,title={Before you move on}]
What does \bn{চড়া} mean? (\textquotedblleft{}To ride, to climb\textquotedblright{}.) Say it once more.
\end{tcolorbox}
